% s-306-p306-m002; original PDF 306, 307
\Aphorism{s-306-p306-m002}{360}
Sittliche Weltordnung — nicht mehr noch minder als eine contradictio in adjecto. Es verlohnte sich einmal, mit kosmischen Flugfahrzeugen und interstellaren Dolmetschern auf die Wanderschaft zu gehen und Umschau zu halten, ob auf sämtlichen bewohnten Planeten eine so curiose Menschheit gediehen ist, die gerade die Begriffe „Welt“ und „Sitte“ zusammenkoppelt, das Engste mit dem Weitesten, das Kirchspiel mit dem Universum, ihre planetarische Privatsache mit der Unendlichkeit. Ein Misanthrop könnte sonst leicht die Vermuthung wagen, dass durch diese Art Transcendenz (zu deutsch: Unverschämtheit) unser tellurisches Völkchen ausgezeichnet, vielleicht gar definirt sei: es hat es an einer reichlichen Lust und Selbstgefälligkeit dabei nicht fehlen lassen. Ein Gott, der über Wechselfälschungen, uneheliche Zusammenkünfte und versäumte Messen Buch führt — derselbe Gott, mit dessen Verabsolutirung zum ens realissimum man sich zwei Jahrtausende lang das Gehirn verrenkt: wer diese schreckliche mésalliance von Göttern einmal nachgedacht hat, wird schwer sich zu entscheiden wissen, was an der Menschheit mehr zu bewundern ist: ihr Wollen oder ihr Können, ihre Fähigkeit zum Unsinn oder ihr Wohlgefallen am Unsinn! — Wie reinlich, wie gut gearbeitet, wie homogen ist dagegen die Lehre von der \Sperr{Wiedergeburt}: beständige Wiederkehr jedes dualistischen Verhältnisses zwischen activer und passiver injuria, mit willkürlicher Vertauschung der beiden Subjecte, des Frevlers und des Leidenden, ewiges unzerstörbares Gleichgewicht zwischen der Ursache als \Sperr{Schuld} und ihrer eigenen Wirkung als \Sperr{Sühne}! Auch hier ist „sittliche Weltordnung“, aber das Sittliche wird von dem inneren Weltmechanismus nebenbei mitbesorgt, nicht ihm als Zweck von aussen gewaltsam aufgedrängt. Und — die Sache ist erkenntnisstheoretisch denkbar, ohne die Grimasse, Verzerrung und Widersprüchlichkeit jenes Gottesbegriffs.
\par

